\section{Conclusiones y trabajo futuro}
\label{sec:conclusiones}

\begin{itemize}
\item KMCPred integra verificación longitudinal de calidad y predicción interpretable por visita sobre una cohorte institucional extensa. La evaluación del modelo de anomalías en producción encontró 83,6 \% de detección de perturbaciones sintéticas con 65,1 alertas/1.000 datos limpios, y la exploración inicial mostró mejoras de AUROC en 13/14 desenlaces frente a una referencia reducida. En la versión publicada se priorizaron cinco modelos clínicos para dos familias de desenlaces a 12 meses: talla baja (AUROC 0,79--0,93) y desarrollo psicomotor bajo (0,77). Estos resultados reflejan validación retrospectiva y no deben trasladarse como rendimiento prospectivo garantizado.
\item Uno de los aprendizajes del proyecto fue que un modelo no basta si el personal no puede entender qué datos analiza o de dónde sale un resultado. Por eso dedicamos parte del desarrollo a organizar los registros con SEVARTI, identificar las visitas disponibles, mostrar las cohortes y separar las funciones clínicas de las técnicas. También fue importante corregir el recuento por niño y advertir cuando una predicción corresponde a una selección anterior. Las gráficas deben leerse siempre junto con el número de niños evaluados, la fecha del corte y el estado del modelo.
\item Como trabajo futuro proponemos empezar por una revisión ciega con especialistas y una evaluación prospectiva sin influir en las decisiones de atención. También necesitamos comprobar el desempeño en otras poblaciones y subgrupos, vigilar los cambios temporales y recalibrar los modelos cuando sea necesario. Quedan pendientes las pruebas de usabilidad, la evaluación de varios usuarios simultáneos y la comparación con reglas clínicas sencillas, como la última talla disponible. Además, sería importante incorporar curvas de crecimiento específicas para niños prematuros antes de las 40 semanas de edad posmenstrual, tomando como referencia los trabajos de Fenton y Kim \cite{r33} y los estándares INTERGROWTH-21st \cite{r34}. Finalmente, los umbrales de alerta deberían acordarse con la Fundación Canguro, teniendo en cuenta tanto los casos identificados como el esfuerzo de revisión que implican.
\end{itemize}
